% =====================================================================
%  Anatomía de un LLM — Fase II — Módulo 11
%  Entrenar redes profundas
%  Compilar: latexmk -r ../latexmkrc M11-redes-profundas.tex
% =====================================================================
\documentclass[11pt,a4paper,twoside,openany]{book}
\usepackage{../estilo-apuntes}

\newcommand{\ApuntesCurso}{Anatomía de un LLM}
\newcommand{\ApuntesModulo}{Módulo 11}

\hypersetup{pdftitle={Módulo 11. Entrenar redes profundas (resumen)},pdfauthor={Anatomía de un LLM}}

% Notación propia del módulo
\newcommand{\T}{^{\top}}
\newcommand{\uno}{\mathbf{1}}
\newcommand{\diag}{\operatorname{diag}}
\newcommand{\Var}{\operatorname{Var}}
\newcommand{\Cov}{\operatorname{Cov}}
\newcommand{\RMS}{\operatorname{RMS}}
\newcommand{\LN}{\operatorname{LN}}
\newcommand{\RMSNorm}{\operatorname{RMSNorm}}
\newcommand{\Norm}{\operatorname{Norm}}
\newcommand{\ReLU}{\operatorname{ReLU}}
\newcommand{\Bern}{\operatorname{Bernoulli}}
\newcommand{\fl}{\operatorname{fl}}
\newcommand{\ulp}{\operatorname{ulp}}
\newcommand{\smin}{\sigma_{\min}}
\newcommand{\Loss}{\mathcal{L}}
\newcommand{\NaN}{\text{\fontseries{m}\ttfamily NaN}}
\newcommand{\Inf}{\text{\fontseries{m}\ttfamily inf}}

\makeatletter
\newlabel{ej:escalar}{{1.3}{2}{Caso escalar}{ejemplo.3}{}}
\newlabel{ej:escalar@cref}{{[ejemplo][3][]1.3}{[1][1][]2}}
\newlabel{ej:potencias}{{1.7}{3}{Potencias de una matriz simétrica}{ejemplo.7}{}}
\newlabel{ej:potencias@cref}{{[ejemplo][7][]1.7}{[1][3][]3}}
\newlabel{ej:pytorch}{{1.12}{4}{Inicialización por defecto de una capa lineal}{ejemplo.12}{}}
\newlabel{ej:pytorch@cref}{{[ejemplo][12][]1.12}{[1][4][]4}}
\newlabel{eje:1-congelada}{{1.1}{5}{}{ejercicio.1.1}{}}
\newlabel{eje:1-congelada@cref}{{[ejercicio][1][1]1.1}{[1][4][]5}}
\newlabel{eje:1-factores}{{1.2}{5}{}{ejercicio.1.2}{}}
\newlabel{eje:1-factores@cref}{{[ejercicio][2][1]1.2}{[1][5][]5}}
\newlabel{eje:1-saturacion}{{1.3}{5}{}{ejercicio.1.3}{}}
\newlabel{eje:1-saturacion@cref}{{[ejercicio][3][1]1.3}{[1][5][]5}}
\newlabel{eje:1-xavier}{{1.4}{5}{}{ejercicio.1.4}{}}
\newlabel{eje:1-xavier@cref}{{[ejercicio][4][1]1.4}{[1][5][]5}}
\newlabel{eje:1-singulares}{{1.5}{5}{}{ejercicio.1.5}{}}
\newlabel{eje:1-singulares@cref}{{[ejercicio][5][1]1.5}{[1][5][]5}}
\newlabel{eje:1-potencias}{{1.6}{5}{}{ejercicio.1.6}{}}
\newlabel{eje:1-potencias@cref}{{[ejercicio][6][1]1.6}{[1][5][]5}}
\newlabel{ej:dos-bloques}{{2.6}{7}{Dos bloques escalares}{ejemplo.6}{}}
\newlabel{ej:dos-bloques@cref}{{[ejemplo][6][]2.6}{[1][7][]7}}
\newlabel{fig:bloque-residual}{{2.1}{8}{Bloque residual. El camino inferior aporta el término $I$ del \cref {teo:jac-residual}; la rama superior aporta $J_\ell $}{figure.2.1}{}}
\newlabel{fig:bloque-residual@cref}{{[figure][1][2]2.1}{[1][7][]8}}
\newlabel{eje:2-anidadas}{{2.1}{8}{}{ejercicio.2.1}{}}
\newlabel{eje:2-anidadas@cref}{{[ejercicio][1][2]2.1}{[1][8][]8}}
\newlabel{eje:2-factor}{{2.2}{9}{}{ejercicio.2.2}{}}
\newlabel{eje:2-factor@cref}{{[ejercicio][2][2]2.2}{[1][9][]9}}
\newlabel{eje:2-matrices}{{2.3}{9}{}{ejercicio.2.3}{}}
\newlabel{eje:2-matrices@cref}{{[ejercicio][3][2]2.3}{[1][9][]9}}
\newlabel{eje:2-flujo}{{2.4}{9}{}{ejercicio.2.4}{}}
\newlabel{eje:2-flujo@cref}{{[ejercicio][4][2]2.4}{[1][9][]9}}
\newlabel{eje:2-cota}{{2.5}{9}{}{ejercicio.2.5}{}}
\newlabel{eje:2-cota@cref}{{[ejercicio][5][2]2.5}{[1][9][]9}}
\newlabel{ej:normas}{{3.4}{11}{}{ejemplo.4}{}}
\newlabel{ej:normas@cref}{{[ejemplo][4][]3.4}{[1][11][]11}}
\newlabel{fig:prenorm}{{3.1}{13}{Bloque decodificador pre-norm, el que se construye en el módulo~14. Cada subbloque normaliza su entrada y suma su salida al flujo residual}{figure.3.1}{}}
\newlabel{fig:prenorm@cref}{{[figure][1][3]3.1}{[1][12][]13}}
\newlabel{eje:3-clasificar}{{3.1}{13}{}{ejercicio.3.1}{}}
\newlabel{eje:3-clasificar@cref}{{[ejercicio][1][3]3.1}{[1][13][]13}}
\newlabel{eje:3-preactivacion}{{3.2}{13}{}{ejercicio.3.2}{}}
\newlabel{eje:3-preactivacion@cref}{{[ejercicio][2][3]3.2}{[1][13][]13}}
\newlabel{eje:3-calculo}{{3.3}{13}{}{ejercicio.3.3}{}}
\newlabel{eje:3-calculo@cref}{{[ejercicio][3][3]3.3}{[1][13][]13}}
\newlabel{eje:3-invariancia}{{3.4}{13}{}{ejercicio.3.4}{}}
\newlabel{eje:3-invariancia@cref}{{[ejercicio][4][3]3.4}{[1][13][]13}}
\newlabel{eje:3-jacobiano}{{3.5}{13}{}{ejercicio.3.5}{}}
\newlabel{eje:3-jacobiano@cref}{{[ejercicio][5][3]3.5}{[1][13][]13}}
\newlabel{eje:3-ortogonal}{{3.6}{13}{}{ejercicio.3.6}{}}
\newlabel{eje:3-ortogonal@cref}{{[ejercicio][6][3]3.6}{[1][13][]13}}
\newlabel{ej:dropout}{{4.4}{15}{}{ejemplo.4}{}}
\newlabel{ej:dropout@cref}{{[ejemplo][4][]4.4}{[1][15][]15}}
\newlabel{ej:recorte}{{4.8}{15}{}{ejemplo.8}{}}
\newlabel{ej:recorte@cref}{{[ejemplo][8][]4.8}{[1][15][]15}}
\newlabel{eje:4-eval}{{4.1}{16}{}{ejercicio.4.1}{}}
\newlabel{eje:4-eval@cref}{{[ejercicio][1][4]4.1}{[1][16][]16}}
\newlabel{eje:4-constante}{{4.2}{16}{}{ejercicio.4.2}{}}
\newlabel{eje:4-constante@cref}{{[ejercicio][2][4]4.2}{[1][16][]16}}
\newlabel{eje:4-mascara}{{4.3}{16}{}{ejercicio.4.3}{}}
\newlabel{eje:4-mascara@cref}{{[ejercicio][3][4]4.3}{[1][16][]16}}
\newlabel{eje:4-recorte}{{4.4}{16}{}{ejercicio.4.4}{}}
\newlabel{eje:4-recorte@cref}{{[ejercicio][4][4]4.4}{[1][16][]16}}
\newlabel{eje:4-global}{{4.5}{16}{}{ejercicio.4.5}{}}
\newlabel{eje:4-global@cref}{{[ejercicio][5][4]4.5}{[1][16][]16}}
\newlabel{eje:4-epsilon}{{4.6}{16}{}{ejercicio.4.6}{}}
\newlabel{eje:4-epsilon@cref}{{[ejercicio][6][4]4.6}{[1][16][]16}}
\newlabel{eje:4-varianza}{{4.7}{17}{}{ejercicio.4.7}{}}
\newlabel{eje:4-varianza@cref}{{[ejercicio][7][4]4.7}{[1][16][]17}}
\newlabel{ej:fp16}{{5.3}{19}{fp16}{ejemplo.3}{}}
\newlabel{ej:fp16@cref}{{[ejemplo][3][]5.3}{[1][18][]19}}
\newlabel{ej:estancamiento}{{5.9}{20}{Estancamiento en fp16 y bf16}{ejemplo.9}{}}
\newlabel{ej:estancamiento@cref}{{[ejemplo][9][]5.9}{[1][20][]20}}
\newlabel{eje:5-riesgo}{{5.1}{21}{}{ejercicio.5.1}{}}
\newlabel{eje:5-riesgo@cref}{{[ejercicio][1][5]5.1}{[1][21][]21}}
\newlabel{eje:5-deducir}{{5.2}{21}{}{ejercicio.5.2}{}}
\newlabel{eje:5-deducir@cref}{{[ejercicio][2][5]5.2}{[1][21][]21}}
\newlabel{eje:5-fp8}{{5.3}{22}{}{ejercicio.5.3}{}}
\newlabel{eje:5-fp8@cref}{{[ejercicio][3][5]5.3}{[1][21][]22}}
\newlabel{eje:5-mil}{{5.4}{22}{}{ejercicio.5.4}{}}
\newlabel{eje:5-mil@cref}{{[ejercicio][4][5]5.4}{[1][22][]22}}
\newlabel{eje:5-rmsnorm}{{5.5}{22}{}{ejercicio.5.5}{}}
\newlabel{eje:5-rmsnorm@cref}{{[ejercicio][5][5]5.5}{[1][22][]22}}
\newlabel{eje:5-absorcion}{{5.6}{22}{}{ejercicio.5.6}{}}
\newlabel{eje:5-absorcion@cref}{{[ejercicio][6][5]5.6}{[1][22][]22}}
\newlabel{eje:5-acumular}{{5.7}{22}{}{ejercicio.5.7}{}}
\newlabel{eje:5-acumular@cref}{{[ejercicio][7][5]5.7}{[1][22][]22}}
\newlabel{ej:lse}{{6.5}{24}{}{ejemplo.5}{}}
\newlabel{ej:lse@cref}{{[ejemplo][5][]6.5}{[1][24][]24}}
\newlabel{ej:log-cero}{{6.7}{24}{}{ejemplo.7}{}}
\newlabel{ej:log-cero@cref}{{[ejemplo][7][]6.7}{[1][24][]24}}
\newlabel{fig:diagnostico}{{6.1}{25}{Árbol de decisión para una pérdida no finita. En todos los casos conviene registrar la norma del gradiente por capa y los histogramas de activaciones}{figure.6.1}{}}
\newlabel{fig:diagnostico@cref}{{[figure][1][6]6.1}{[1][25][]25}}
\newlabel{eje:6-diagnostico}{{6.1}{26}{}{ejercicio.6.1}{}}
\newlabel{eje:6-diagnostico@cref}{{[ejercicio][1][6]6.1}{[1][26][]26}}
\newlabel{eje:6-lse}{{6.2}{26}{}{ejercicio.6.2}{}}
\newlabel{eje:6-lse@cref}{{[ejercicio][2][6]6.2}{[1][26][]26}}
\newlabel{eje:6-softmax}{{6.3}{26}{}{ejercicio.6.3}{}}
\newlabel{eje:6-softmax@cref}{{[ejercicio][3][6]6.3}{[1][26][]26}}
\newlabel{eje:6-rango}{{6.4}{26}{}{ejercicio.6.4}{}}
\newlabel{eje:6-rango@cref}{{[ejercicio][4][6]6.4}{[1][26][]26}}
\newlabel{eje:6-ce}{{6.5}{26}{}{ejercicio.6.5}{}}
\newlabel{eje:6-ce@cref}{{[ejercicio][5][6]6.5}{[1][26][]26}}
\newlabel{eje:6-cotas}{{6.6}{26}{}{ejercicio.6.6}{}}
\newlabel{eje:6-cotas@cref}{{[ejercicio][6][6]6.6}{[1][26][]26}}
\newlabel{eje:6-convexa}{{6.7}{26}{}{ejercicio.6.7}{}}
\newlabel{eje:6-convexa@cref}{{[ejercicio][7][6]6.7}{[1][26][]26}}
\newlabel{ap:soluciones}{{A}{27}{Soluciones de los ejercicios}{appendix.A}{}}
\newlabel{ap:soluciones@cref}{{[appendix][1][2147483647]A}{[1][27][]27}}
\newlabel{sol:eje:1-congelada}{{A.1}{27}{Tema 1}{section*.6}{}}
\newlabel{sol:eje:1-congelada@cref}{{[subappendix][1][2147483647,1]A.1}{[1][27][]27}}
\newlabel{sol:eje:1-factores}{{A.1}{27}{Tema 1}{section*.7}{}}
\newlabel{sol:eje:1-factores@cref}{{[subappendix][1][2147483647,1]A.1}{[1][27][]27}}
\newlabel{sol:eje:1-saturacion}{{A.1}{28}{Tema 1}{section*.8}{}}
\newlabel{sol:eje:1-saturacion@cref}{{[subappendix][1][2147483647,1]A.1}{[1][28][]28}}
\newlabel{sol:eje:1-xavier}{{A.1}{28}{Tema 1}{section*.9}{}}
\newlabel{sol:eje:1-xavier@cref}{{[subappendix][1][2147483647,1]A.1}{[1][28][]28}}
\newlabel{sol:eje:1-singulares}{{A.1}{28}{Tema 1}{section*.10}{}}
\newlabel{sol:eje:1-singulares@cref}{{[subappendix][1][2147483647,1]A.1}{[1][28][]28}}
\newlabel{sol:eje:1-potencias}{{A.1}{29}{Tema 1}{section*.11}{}}
\newlabel{sol:eje:1-potencias@cref}{{[subappendix][1][2147483647,1]A.1}{[1][29][]29}}
\newlabel{sol:eje:2-anidadas}{{A.2}{29}{Tema 2}{section*.12}{}}
\newlabel{sol:eje:2-anidadas@cref}{{[subappendix][2][2147483647,1]A.2}{[1][29][]29}}
\newlabel{sol:eje:2-factor}{{A.2}{30}{Tema 2}{section*.13}{}}
\newlabel{sol:eje:2-factor@cref}{{[subappendix][2][2147483647,1]A.2}{[1][30][]30}}
\newlabel{sol:eje:2-matrices}{{A.2}{30}{Tema 2}{section*.14}{}}
\newlabel{sol:eje:2-matrices@cref}{{[subappendix][2][2147483647,1]A.2}{[1][30][]30}}
\newlabel{sol:eje:2-flujo}{{A.2}{31}{Tema 2}{section*.15}{}}
\newlabel{sol:eje:2-flujo@cref}{{[subappendix][2][2147483647,1]A.2}{[1][31][]31}}
\newlabel{sol:eje:2-cota}{{A.2}{31}{Tema 2}{section*.16}{}}
\newlabel{sol:eje:2-cota@cref}{{[subappendix][2][2147483647,1]A.2}{[1][31][]31}}
\newlabel{sol:eje:3-clasificar}{{A.3}{31}{Tema 3}{section*.17}{}}
\newlabel{sol:eje:3-clasificar@cref}{{[subappendix][3][2147483647,1]A.3}{[1][31][]31}}
\newlabel{sol:eje:3-preactivacion}{{A.3}{32}{Tema 3}{section*.18}{}}
\newlabel{sol:eje:3-preactivacion@cref}{{[subappendix][3][2147483647,1]A.3}{[1][32][]32}}
\newlabel{sol:eje:3-calculo}{{A.3}{32}{Tema 3}{section*.19}{}}
\newlabel{sol:eje:3-calculo@cref}{{[subappendix][3][2147483647,1]A.3}{[1][32][]32}}
\newlabel{sol:eje:3-invariancia}{{A.3}{33}{Tema 3}{section*.20}{}}
\newlabel{sol:eje:3-invariancia@cref}{{[subappendix][3][2147483647,1]A.3}{[1][33][]33}}
\newlabel{sol:eje:3-jacobiano}{{A.3}{33}{Tema 3}{section*.21}{}}
\newlabel{sol:eje:3-jacobiano@cref}{{[subappendix][3][2147483647,1]A.3}{[1][33][]33}}
\newlabel{sol:eje:3-ortogonal}{{A.3}{33}{Tema 3}{section*.22}{}}
\newlabel{sol:eje:3-ortogonal@cref}{{[subappendix][3][2147483647,1]A.3}{[1][33][]33}}
\newlabel{sol:eje:4-eval}{{A.4}{34}{Tema 4}{section*.23}{}}
\newlabel{sol:eje:4-eval@cref}{{[subappendix][4][2147483647,1]A.4}{[1][34][]34}}
\newlabel{sol:eje:4-constante}{{A.4}{34}{Tema 4}{section*.24}{}}
\newlabel{sol:eje:4-constante@cref}{{[subappendix][4][2147483647,1]A.4}{[1][34][]34}}
\newlabel{sol:eje:4-mascara}{{A.4}{35}{Tema 4}{section*.25}{}}
\newlabel{sol:eje:4-mascara@cref}{{[subappendix][4][2147483647,1]A.4}{[1][34][]35}}
\newlabel{sol:eje:4-recorte}{{A.4}{35}{Tema 4}{section*.26}{}}
\newlabel{sol:eje:4-recorte@cref}{{[subappendix][4][2147483647,1]A.4}{[1][35][]35}}
\newlabel{sol:eje:4-global}{{A.4}{35}{Tema 4}{section*.27}{}}
\newlabel{sol:eje:4-global@cref}{{[subappendix][4][2147483647,1]A.4}{[1][35][]35}}
\newlabel{sol:eje:4-epsilon}{{A.4}{36}{Tema 4}{section*.28}{}}
\newlabel{sol:eje:4-epsilon@cref}{{[subappendix][4][2147483647,1]A.4}{[1][36][]36}}
\newlabel{sol:eje:4-varianza}{{A.4}{36}{Tema 4}{section*.29}{}}
\newlabel{sol:eje:4-varianza@cref}{{[subappendix][4][2147483647,1]A.4}{[1][36][]36}}
\newlabel{sol:eje:5-riesgo}{{A.5}{37}{Tema 5}{section*.30}{}}
\newlabel{sol:eje:5-riesgo@cref}{{[subappendix][5][2147483647,1]A.5}{[1][36][]37}}
\newlabel{sol:eje:5-deducir}{{A.5}{37}{Tema 5}{section*.31}{}}
\newlabel{sol:eje:5-deducir@cref}{{[subappendix][5][2147483647,1]A.5}{[1][37][]37}}
\newlabel{sol:eje:5-fp8}{{A.5}{37}{Tema 5}{section*.32}{}}
\newlabel{sol:eje:5-fp8@cref}{{[subappendix][5][2147483647,1]A.5}{[1][37][]37}}
\newlabel{sol:eje:5-mil}{{A.5}{38}{Tema 5}{section*.33}{}}
\newlabel{sol:eje:5-mil@cref}{{[subappendix][5][2147483647,1]A.5}{[1][38][]38}}
\newlabel{sol:eje:5-rmsnorm}{{A.5}{38}{Tema 5}{section*.34}{}}
\newlabel{sol:eje:5-rmsnorm@cref}{{[subappendix][5][2147483647,1]A.5}{[1][38][]38}}
\newlabel{sol:eje:5-absorcion}{{A.5}{39}{Tema 5}{section*.35}{}}
\newlabel{sol:eje:5-absorcion@cref}{{[subappendix][5][2147483647,1]A.5}{[1][39][]39}}
\newlabel{sol:eje:5-acumular}{{A.5}{39}{Tema 5}{section*.36}{}}
\newlabel{sol:eje:5-acumular@cref}{{[subappendix][5][2147483647,1]A.5}{[1][39][]39}}
\newlabel{sol:eje:6-diagnostico}{{A.6}{40}{Tema 6}{section*.37}{}}
\newlabel{sol:eje:6-diagnostico@cref}{{[subappendix][6][2147483647,1]A.6}{[1][40][]40}}
\newlabel{sol:eje:6-lse}{{A.6}{40}{Tema 6}{section*.38}{}}
\newlabel{sol:eje:6-lse@cref}{{[subappendix][6][2147483647,1]A.6}{[1][40][]40}}
\newlabel{sol:eje:6-softmax}{{A.6}{40}{Tema 6}{section*.39}{}}
\newlabel{sol:eje:6-softmax@cref}{{[subappendix][6][2147483647,1]A.6}{[1][40][]40}}
\newlabel{sol:eje:6-rango}{{A.6}{41}{Tema 6}{section*.40}{}}
\newlabel{sol:eje:6-rango@cref}{{[subappendix][6][2147483647,1]A.6}{[1][41][]41}}
\newlabel{sol:eje:6-ce}{{A.6}{41}{Tema 6}{section*.41}{}}
\newlabel{sol:eje:6-ce@cref}{{[subappendix][6][2147483647,1]A.6}{[1][41][]41}}
\newlabel{sol:eje:6-cotas}{{A.6}{41}{Tema 6}{section*.42}{}}
\newlabel{sol:eje:6-cotas@cref}{{[subappendix][6][2147483647,1]A.6}{[1][41][]41}}
\newlabel{sol:eje:6-convexa}{{A.6}{42}{Tema 6}{section*.43}{}}
\newlabel{sol:eje:6-convexa@cref}{{[subappendix][6][2147483647,1]A.6}{[1][42][]42}}
\makeatother

\begin{document}
\frontmatter
\begin{titlepage}
\centering
\vspace*{2.5cm}
{\large\scshape Anatomía de un LLM}\par
\vspace{0.4cm}
{\normalsize Fundamentos matemáticos de los modelos de lenguaje}\par
\vspace{2.2cm}
\rule{\textwidth}{0.8pt}\par\vspace{0.5cm}
{\small\scshape Fase II \textperiodcentered\ Aprendizaje automático y redes neuronales}\par\vspace{0.3cm}
{\Huge\bfseries Módulo 11\par}\vspace{0.3cm}
{\LARGE Entrenar redes profundas\par}
\vspace{0.5cm}\rule{\textwidth}{0.4pt}\par
\vspace{1.6cm}
{\large Resumen del módulo}\par
\vfill
{\small Edición de octubre de 2026}\par
\end{titlepage}

\chapter*{Sobre este resumen}
\addcontentsline{toc}{chapter}{Sobre este resumen}
\markboth{Sobre este resumen}{}

Este documento reúne lo indispensable de los apuntes del Módulo 11: el texto expositivo, todas las definiciones y todos los resultados (teoremas, proposiciones, lemas y corolarios), las tablas y las observaciones de contenido. Se han suprimido las demostraciones, los ejemplos, las lecturas, las figuras, los ejercicios y sus soluciones. La numeración de temas, secciones, enunciados y ecuaciones es la misma que en los apuntes, de modo que cualquier referencia puede buscarse allí; las que remiten a elementos suprimidos (un ejemplo, una figura, una ecuación de una demostración) conservan su número en los apuntes.

\section*{Notación}
\begin{center}
\begin{tabular}{@{}ll@{}}
\toprule
Símbolo & Significado \\
\midrule
$h_\ell\in\R^{d}$, $z_\ell=W_\ell h_{\ell-1}$ & activación y preactivación de la capa $\ell$, con $\ell=1,\dots,L$ \\
$J_\ell=\partial h_\ell/\partial h_{\ell-1}$ & jacobiano de la capa (o de la rama residual, tema~2) \\
$\Loss$, $g=\nabla_\theta\Loss$ & pérdida y gradiente respecto de los parámetros \\
$\norm{A}_2$, $\smin(A)$ & norma espectral (mayor valor singular) y menor valor singular \\
$\Var$, $\E$ & varianza y esperanza \\
$\mu$, $\sigma^{2}$, $\RMS(x)$ & media, varianza y raíz cuadrática media de un vector \\
$\gamma,\beta$; $g$ & ganancia y sesgo de LayerNorm; ganancia de RMSNorm \\
$\epsilon$ & constante de estabilidad dentro de una raíz \\
$p$, $m_i\sim\Bern(1-p)$ & probabilidad de apagado del dropout y máscara \\
$c$ & umbral del recorte del gradiente \\
$e$, $m$ & bits de exponente y de mantisa de un formato de coma flotante \\
$\varepsilon=2^{-m}$, $u=\varepsilon/2$ & épsilon de máquina y unidad de redondeo \\
$\fl(x)$ & número representable más cercano a $x$ \\
$\LSE(z)$ & log-sum-exp, $\ln\sum_i e^{z_i}$ \\
$u\odot w$, $\uno$ & producto elemento a elemento; vector de unos \\
\bottomrule
\end{tabular}
\end{center}


\tableofcontents

\mainmatter
% =====================================================================
\chapter{Gradientes que se desvanecen o explotan}
\label{tema:1}
% =====================================================================

En una red profunda, el gradiente que recibe una capa es el producto de los jacobianos de todas las capas que tiene por encima. Este tema formaliza esa observación, acota la norma del producto mediante valores singulares y deduce de ella los dos fenómenos que dificultan el entrenamiento en profundidad: el desvanecimiento y la explosión del gradiente. Se estudian después sus dos causas principales, la escala de los pesos y la saturación de las activaciones, y se obtiene la inicialización que mantiene constante la varianza de la señal. Se sigue~\cite[§8.2.5, §10.7]{gbc}, \cite[§5.4 y §9.7]{d2l} y los artículos de Pascanu, Mikolov y Bengio~\cite{pascanu}, Glorot y Bengio~\cite{glorot} y He \emph{et al.}~\cite{he-init}.

\section{El gradiente como producto de jacobianos}

\begin{definicion}[Red profunda de anchura constante]\label{def:red}
Sea $L\ge 1$. Una \emph{red de $L$ capas} con anchura $d$ y activación $f\colon\R\to\R$ derivable transforma una entrada $h_0\in\R^{d}$ mediante
\[
z_\ell=W_\ell h_{\ell-1},\qquad h_\ell=f(z_\ell),\qquad \ell=1,\dots,L,
\]
con $W_\ell\in\R^{d\times d}$ y $f$ aplicada componente a componente. La pérdida $\Loss$ depende de la entrada solo a través de $h_L$. El \emph{jacobiano de la capa $\ell$} es
\[
J_\ell=\frac{\partial h_\ell}{\partial h_{\ell-1}}=\diag\bigl(f'(z_\ell)\bigr)\,W_\ell .
\]
\end{definicion}

Los sesgos se omiten porque no intervienen en los jacobianos respecto de las activaciones. La anchura constante simplifica la notación; todos los resultados se extienden a anchuras variables.

\begin{proposicion}[Gradiente en profundidad]\label{prop:grad-prof}
Con la notación de la \cref{def:red}, escribiendo los gradientes como vectores columna,
\[
\nabla_{h_{\ell}}\Loss=J_{\ell+1}\T J_{\ell+2}\T\cdots J_L\T\,\nabla_{h_L}\Loss,
\qquad
\frac{\partial\Loss}{\partial W_\ell}=\Bigl(\diag\bigl(f'(z_\ell)\bigr)\nabla_{h_\ell}\Loss\Bigr)\,h_{\ell-1}\T .
\]
\end{proposicion}

El gradiente de los pesos de la capa $\ell$ contiene, por tanto, el producto de los $L-\ell$ jacobianos que tiene encima. Las primeras capas atraviesan casi todos.

\addtocounter{definicion}{1}%

\begin{table}[h]
\centering
\begin{tabular}{@{}cccl@{}}
\toprule
Factor por capa $a$ & $L=10$ & $L=24$ & Efecto \\
\midrule
$0{,}25$ & $9{,}5\cdot 10^{-7}$ & $3{,}6\cdot 10^{-15}$ & desvanecimiento extremo \\
$0{,}5$ & $9{,}8\cdot 10^{-4}$ & $6{,}0\cdot 10^{-8}$ & las primeras capas no aprenden \\
$1$ & $1$ & $1$ & caso ideal \\
$1{,}1$ & $2{,}59$ & $9{,}85$ & crecimiento moderado \\
$1{,}5$ & $57{,}7$ & $1{,}7\cdot 10^{4}$ & explosión \\
\bottomrule
\end{tabular}
\caption{Factor $a^{L}$ por el que se multiplica el gradiente al atravesar $L$ capas escalares.}
\label{tab:factores}
\end{table}

\section{Cotas mediante valores singulares}

Para matrices, el papel del factor $a$ lo desempeñan los valores singulares. Para $A\in\R^{d\times d}$ se escribe $\norm{A}_2=\max_{\norm{x}=1}\norm{Ax}$ (el mayor valor singular, módulo~3) y $\smin(A)=\min_{\norm{x}=1}\norm{Ax}$ (el menor). Ambos coinciden para $A$ y $A\T$.

\begin{lema}\label{lem:producto}
Para $A,B\in\R^{d\times d}$ se cumple $\smin(A)\,\smin(B)\le\smin(AB)\le\norm{AB}_2\le\norm{A}_2\norm{B}_2$.
\end{lema}

\begin{teorema}[Cotas del gradiente en profundidad]\label{teo:cotas}
Con la notación de la \cref{def:red},
\[
\prod_{\ell=1}^{L}\smin(J_\ell)\;\le\;\frac{\norm{\nabla_{h_0}\Loss}}{\norm{\nabla_{h_L}\Loss}}\;\le\;\prod_{\ell=1}^{L}\norm{J_\ell}_2\;\le\;\prod_{\ell=1}^{L}\gamma\,\norm{W_\ell}_2,
\qquad \gamma=\sup_{t\in\R}\abs{f'(t)} .
\]
\end{teorema}

\begin{corolario}[Condición suficiente de desvanecimiento]\label{cor:pascanu}
Si existe $\lambda<1$ tal que $\gamma\norm{W_\ell}_2\le\lambda$ para todo $\ell$, entonces $\norm{\nabla_{h_0}\Loss}\le\lambda^{L}\norm{\nabla_{h_L}\Loss}$, que tiende a~0 exponencialmente con $L$. En consecuencia, que el gradiente explote exige que $\norm{W_\ell}_2>1/\gamma$ para alguna capa.
\end{corolario}

En una red recurrente la misma matriz $W$ se aplica una vez por paso de tiempo, de modo que el corolario se lee con $L$ igual a la longitud de la secuencia. Es el resultado de~\cite[§2.1]{pascanu}: con $\gamma=1$ para $\tanh$ y $\gamma=1/4$ para la sigmoide logística, basta $\norm{W}_2<1/\gamma$ para que el gradiente se desvanezca a lo largo del tiempo. Goodfellow \emph{et al.}~\cite[§10.7]{gbc} y Zhang \emph{et al.}~\cite[§9.7.2]{d2l} presentan el caso lineal, en el que el gradiente contiene las potencias $W^{t}$.

\addtocounter{definicion}{1}%

\begin{observacion}
Los productos de muchos factores grandes crean en la superficie de pérdida regiones casi verticales, llamadas \emph{acantilados}~\cite[§8.2.4]{gbc}. Un paso de gradiente en una de ellas puede expulsar los parámetros lejos de la región que se estaba explorando. El recorte del gradiente del \cref{tema:4} es el remedio estándar.
\end{observacion}

\section{Saturación de las activaciones}

El factor $\diag(f'(z_\ell))$ del jacobiano depende de las preactivaciones. Para las activaciones del módulo~9:
\begin{itemize}[leftmargin=1.4em]
  \item $\tanh'(t)=1-\tanh^{2}(t)\in(0,1]$, con máximo en $t=0$. Vale $0{,}420$ en $t=1$, $0{,}0707$ en $t=2$ y $0{,}0099$ en $t=3$.
  \item La sigmoide $s(t)=1/(1+e^{-t})$ cumple $s'(t)=s(t)(1-s(t))\le 1/4$.
  \item $\ReLU'(t)=1$ si $t>0$ y $0$ si $t<0$.
\end{itemize}

Una unidad con $\abs{z}$ grande está \emph{saturada}: su derivada es casi nula y anula la fila correspondiente del jacobiano. Si una unidad ReLU tiene preactivación negativa para todas las entradas, su gradiente es exactamente cero y deja de aprender. Por el \cref{teo:cotas}, basta que una sola capa tenga $\smin(J_\ell)$ pequeño para que el producto completo pueda serlo.

\section{Propagación de la varianza e inicialización}

La escala de los pesos se elige para que la varianza de la señal no crezca ni decrezca de una capa a la siguiente. El resultado clave es el siguiente~\cite[§5.4.2]{d2l}.

\begin{teorema}[Propagación de la varianza]\label{teo:varianza}
Sea $y=\sum_{i=1}^{n}w_ix_i$, donde las $2n$ variables $w_1,\dots,w_n,x_1,\dots,x_n$ son independientes, $\E[w_i]=0$, $\Var(w_i)=\sigma^{2}$ y $\E[x_i^{2}]=q$ para todo $i$. Entonces $\E[y]=0$ y
\[
\Var(y)=n\,\sigma^{2}\,q .
\]
Si además $\E[x_i]=0$, entonces $q=\Var(x_i)$.
\end{teorema}

\begin{corolario}[Inicialización de Xavier]\label{cor:xavier}
En una red lineal ($f(t)=t$) con entradas centradas, $\Var(z_\ell)=(n\sigma^{2})^{\ell}\Var(h_0)$. La varianza se conserva en todas las capas si y solo si $\sigma^{2}=1/n$. La pasada hacia atrás multiplica por las matrices traspuestas, cuyo número de entradas por fila es el de salidas, y la varianza del gradiente se conserva con $\sigma^{2}=1/n_{\text{sal}}$; Glorot y Bengio~\cite{glorot} proponen el compromiso $\sigma^{2}=2/(n_{\text{ent}}+n_{\text{sal}})$.
\end{corolario}

\begin{proposicion}[Inicialización de He]\label{prop:he}
Sea $z$ una variable aleatoria con densidad simétrica respecto de~0 y varianza $v$. Entonces $\E[\ReLU(z)^{2}]=v/2$. En consecuencia, en una red ReLU la varianza de las preactivaciones se conserva con $\sigma^{2}=2/n$~\cite{he-init}.
\end{proposicion}

\addtocounter{definicion}{1}%

\addtocounter{definicion}{1}%


% =====================================================================
\chapter{Conexiones residuales}
\label{tema:2}
% =====================================================================

Una conexión residual suma a la salida de un bloque su propia entrada. Con este cambio mínimo, el jacobiano de cada bloque pasa a ser la identidad más una corrección, y el producto de jacobianos del \cref{tema:1} contiene siempre un término que no se atenúa con la profundidad. Este tema define las redes residuales, demuestra sus propiedades de propagación hacia delante y hacia atrás, y describe el flujo residual de un Transformer. Se sigue el artículo original de He \emph{et al.}~\cite{resnet}, su continuación sobre aplicaciones identidad~\cite{resnet-id}, \cite[§8.6]{d2l} y \cite[§9.5]{bishopdl}.

\section{Clases de funciones anidadas}

Zhang \emph{et al.}~\cite[§8.6.1]{d2l} motivan las redes residuales con un argumento sobre clases de funciones. Sea $\mathcal{F}_L$ el conjunto de funciones que puede representar una arquitectura de $L$ bloques. Al añadir un bloque se esperaría $\mathcal{F}_L\subseteq\mathcal{F}_{L+1}$: una red más profunda debería poder imitar a la menos profunda. En una red plana esto exige que el bloque nuevo aprenda exactamente la identidad, lo que no es fácil con capas no lineales; empíricamente, las redes planas muy profundas obtienen \emph{mayor} error de entrenamiento que otras menos profundas~\cite{resnet}.

\begin{definicion}[Bloque y red residual]\label{def:residual}
Un \emph{bloque residual} es una aplicación $\R^{d}\to\R^{d}$ de la forma
\[
h\longmapsto h+F(h;\theta),
\]
donde $F$, la \emph{rama residual}, es una red diferenciable con parámetros $\theta$. Una \emph{red residual} de $L$ bloques calcula $h_\ell=h_{\ell-1}+F_\ell(h_{\ell-1})$ para $\ell=1,\dots,L$.
\end{definicion}

\begin{proposicion}[Anidamiento]\label{prop:anidamiento}
Si cada rama $F_\ell$ termina en una capa lineal $u\mapsto Vu$, la red residual de $L+1$ bloques puede representar cualquier función que represente la de $L$ bloques.
\end{proposicion}

La proposición explica también la idea de \emph{residuo}: el bloque no tiene que aprender la transformación deseada $H(h)$, sino la diferencia $F(h)=H(h)-h$, que es pequeña cuando $H$ se parece a la identidad.

\section{Propagación en una red residual}

\begin{proposicion}[Propagación hacia delante]\label{prop:adelante}
En una red residual, para $0\le\ell<L$,
\[
h_L=h_\ell+\sum_{k=\ell+1}^{L}F_k(h_{k-1}).
\]
\end{proposicion}

La salida es la entrada de cualquier bloque más la suma de las aportaciones de los bloques posteriores. Esta lectura aditiva, destacada en~\cite[§2]{resnet-id}, contrasta con la composición $f_L\circ\cdots\circ f_1$ de una red plana.

\begin{teorema}[Jacobiano de una red residual]\label{teo:jac-residual}
Sea $J_\ell=\partial F_\ell/\partial h_{\ell-1}$, evaluado en $h_{\ell-1}$. Entonces
\[
\frac{\partial h_L}{\partial h_0}=(I+J_L)(I+J_{L-1})\cdots(I+J_1)
=I+\sum_{\ell=1}^{L}J_\ell+\sum_{1\le\ell<k\le L}J_kJ_\ell+\cdots+J_L\cdots J_1,
\]
donde el desarrollo tiene $2^{L}$ términos, uno por cada subconjunto de bloques, con los factores en orden decreciente de índice. En particular, contiene el término $I$.
\end{teorema}

\begin{corolario}\label{cor:grad-residual}
$\nabla_{h_0}\Loss=\nabla_{h_L}\Loss+\Bigl(\sum_{\ell}J_\ell+\sum_{\ell<k}J_kJ_\ell+\cdots\Bigr)\T\nabla_{h_L}\Loss$. El gradiente de la salida llega a la entrada sin atenuar, sumado a correcciones.
\end{corolario}

Cada término del desarrollo corresponde a un camino por la red que atraviesa un subconjunto de ramas y salta las demás. Una red plana solo dispone del camino $J_L\cdots J_1$, que es el último término.

\addtocounter{definicion}{1}%

\begin{proposicion}[Cota inferior]\label{prop:cota-residual}
Si $\norm{J_\ell}_2\le a<1$ para todo $\ell$, entonces
\[
\smin\Bigl(\frac{\partial h_L}{\partial h_0}\Bigr)\ge(1-a)^{L}>0 .
\]
En particular, ningún vector de gradiente no nulo en la salida se anula en la entrada.
\end{proposicion}

En una red plana, en cambio, basta un $J_\ell$ con núcleo no trivial (una capa ReLU con unidades apagadas, por ejemplo) para que algún gradiente se anule. La cota de la \cref{prop:cota-residual} no impide el desvanecimiento si $a$ es grande, pero al inicio del entrenamiento las ramas tienen jacobianos pequeños y $(1\pm a)^{L}$ queda cerca de~1: con $a=0{,}01$ y $L=24$, entre $0{,}99^{24}=0{,}786$ y $1{,}01^{24}=1{,}27$.

\addtocounter{figure}{1}%

\section{El flujo residual de un Transformer}

En un Transformer, el vector $h_\ell\in\R^{d}$ asociado a cada token recorre todos los bloques. Cada subbloque de atención o de perceptrón multicapa lee de él y le suma su aportación; por eso se le llama \emph{flujo residual}. La \cref{prop:adelante} dice que el flujo de la última capa es la suma del embedding inicial y de todas las aportaciones. En el módulo~23 se usará esta descomposición para interpretar qué escribe cada bloque.

\begin{proposicion}[Crecimiento del flujo residual]\label{prop:crecimiento}
Supóngase que, al inicio del entrenamiento, las aportaciones $F_k(h_{k-1})$ de una componente del flujo son variables aleatorias de media nula, incorreladas entre sí y con $h_0$, y con varianza $v$ cada una. Entonces $\Var\bigl((h_L)_i\bigr)=\Var\bigl((h_0)_i\bigr)+Lv$: la desviación típica del flujo crece como $\sqrt{L}$.
\end{proposicion}

Este crecimiento motiva dos prácticas de los modelos de lenguaje. GPT-2~\cite[§2.3]{gpt2} multiplica en la inicialización los pesos de las capas que escriben en el flujo por $1/\sqrt{N}$, con $N$ el número de capas residuales, lo que reduce cada varianza $v$ en un factor $N$. Y los Transformer pre-norm añaden una normalización final antes de la capa de salida (\cref{tema:3}).

\addtocounter{definicion}{1}%


% =====================================================================
\chapter{Normalización: LayerNorm, RMSNorm y pre-norm}
\label{tema:3}
% =====================================================================

Normalizar una capa consiste en reescalar cada vector de activaciones con estadísticos calculados sobre él mismo, de modo que su escala no derive de una capa a otra. Este tema define la normalización por lotes, la normalización por capa (LayerNorm) y su simplificación RMSNorm, demuestra sus propiedades de invariancia, calcula el jacobiano de RMSNorm y analiza dónde debe colocarse la normalización respecto de la conexión residual. Se sigue~\cite[§7.4]{bishopdl}, \cite[§8.5]{d2l} y \cite[§8.7.1]{gbc}, junto con los artículos de Ioffe y Szegedy~\cite{ioffe}, Ba, Kiros y Hinton~\cite{layernorm}, Zhang y Sennrich~\cite{rmsnorm} y Xiong \emph{et al.}~\cite{xiong}.

\section{Normalización por lotes y por capa}

\begin{definicion}[Normalización por lotes]\label{def:batchnorm}
Sea un minibatch de $B$ vectores $x^{(1)},\dots,x^{(B)}\in\R^{d}$. La \emph{normalización por lotes} (BatchNorm) normaliza cada \emph{componente} $i$ con la media y la varianza de esa componente sobre el lote:
\[
\hat x^{(b)}_i=\frac{x^{(b)}_i-\mu_i}{\sqrt{\sigma_i^{2}+\epsilon}},\qquad
\mu_i=\frac1B\sum_{b=1}^{B}x^{(b)}_i,\quad \sigma_i^{2}=\frac1B\sum_{b=1}^{B}\bigl(x^{(b)}_i-\mu_i\bigr)^{2},
\]
y después aplica una ganancia y un sesgo aprendidos por componente~\cite{ioffe}.
\end{definicion}

La salida de BatchNorm para un ejemplo depende de los demás ejemplos del lote; en inferencia se sustituyen $\mu_i$ y $\sigma_i^{2}$ por medias móviles estimadas durante el entrenamiento~\cite[§8.5.2]{d2l}. Para secuencias de longitud variable y para la generación token a token, esta dependencia es un inconveniente. La normalización por capa la elimina normalizando sobre las componentes de cada vector.

\begin{definicion}[LayerNorm]\label{def:layernorm}
Sea $x\in\R^{d}$, y sean $\gamma,\beta\in\R^{d}$ parámetros y $\epsilon\ge 0$. Con
\[
\mu=\frac1d\sum_{i=1}^{d}x_i,\qquad \sigma^{2}=\frac1d\sum_{i=1}^{d}(x_i-\mu)^{2},
\]
la \emph{normalización por capa} es
\[
\LN(x)_i=\gamma_i\,\frac{x_i-\mu}{\sqrt{\sigma^{2}+\epsilon}}+\beta_i,\qquad i=1,\dots,d~\cite{layernorm}.
\]
\end{definicion}

\begin{definicion}[RMSNorm]\label{def:rmsnorm}
Sea $x\in\R^{d}$, $g\in\R^{d}$ y $\epsilon\ge 0$. Con $\RMS(x)=\sqrt{\frac1d\sum_{i=1}^{d}x_i^{2}+\epsilon}$, se define
\[
\RMSNorm(x)_i=g_i\,\frac{x_i}{\RMS(x)},\qquad i=1,\dots,d~\cite{rmsnorm}.
\]
\end{definicion}

RMSNorm suprime el centrado y el sesgo $\beta$: solo fija la escala. En un Transformer ambas se aplican por separado al vector de cada token, por lo que funcionan igual en entrenamiento y en inferencia y con lotes de cualquier tamaño.

\addtocounter{definicion}{1}%

\begin{proposicion}[Estadísticos de la salida]\label{prop:estad-salida}
Sea $\epsilon=0$, $\gamma=g=\uno$ y $\beta=0$.
\begin{enumerate}
  \item Si $x$ no es constante, $y=\LN(x)$ cumple $\frac1d\sum_iy_i=0$ y $\frac1d\sum_iy_i^{2}=1$.
  \item Si $x\ne 0$, $y=\RMSNorm(x)$ cumple $\frac1d\sum_iy_i^{2}=1$, es decir, $\norm{y}=\sqrt d$.
\end{enumerate}
\end{proposicion}

\section{Invariancias}

\begin{proposicion}[Invariancias]\label{prop:invariancias}
Sea $\epsilon=0$, $\alpha>0$ y $c\in\R$.
\begin{enumerate}
  \item $\LN(\alpha x+c\uno)=\LN(x)$ para todo $x$ no constante.
  \item $\RMSNorm(\alpha x)=\RMSNorm(x)$ para todo $x\ne 0$, pero en general $\RMSNorm(x+c\uno)\ne\RMSNorm(x)$.
  \item Si la entrada de la normalización es $x=Wu$ con $W\in\R^{d\times n}$, la salida no cambia al sustituir $W$ por $\alpha W$, para LayerNorm y para RMSNorm.
\end{enumerate}
\end{proposicion}

La propiedad (iii) figura en la tabla de invariancias de Ba \emph{et al.}~\cite[§5.1]{layernorm}: la salida no depende de la escala de los pesos de la capa anterior. Zhang y Sennrich~\cite{rmsnorm} conservan esta invariancia de escala, que consideran la responsable del buen comportamiento de LayerNorm, y renuncian a la de traslación, con lo que se ahorra el cálculo de la media.

\begin{teorema}[Jacobiano de RMSNorm]\label{teo:jac-rms}
Sea $\epsilon=0$, $x\ne 0$ y $r=\RMS(x)$. El jacobiano de $y=\RMSNorm(x)$ es
\[
J=\frac{\partial y}{\partial x}=\frac1r\diag(g)\Bigl(I-\frac{xx\T}{d\,r^{2}}\Bigr)
=\frac1r\diag(g)\Bigl(I-\frac{xx\T}{\norm{x}^{2}}\Bigr).
\]
En consecuencia, $Jx=0$ y, para todo $v\in\R^{d}$, $x\T(J\T v)=0$: el gradiente que atraviesa RMSNorm hacia atrás es ortogonal a $x$, y su norma es a lo sumo $\max_i\abs{g_i}\,\norm{v}/r$.
\end{teorema}

La identidad $Jx=0$ es la versión infinitesimal de la invariancia de escala: al derivar $\RMSNorm((1+t)x)=\RMSNorm(x)$ en $t=0$ se obtiene $Jx=0$. La cota $\norm{v}/r$ muestra el efecto estabilizador: si una capa anterior infla sus activaciones, $r$ crece y el gradiente que se le devuelve se reduce en la misma proporción.

\begin{observacion}
Un cálculo análogo, con $\hat y=(x-\mu\uno)/\sigma$, da para LayerNorm con $\gamma=\uno$, $\beta=0$ y $\epsilon=0$
\[
\frac{\partial\hat y}{\partial x}=\frac1\sigma\Bigl(I-\frac{\uno\uno\T}{d}-\frac{\hat y\hat y\T}{d}\Bigr),
\]
que anula tanto $\uno$ como $x-\mu\uno$, de acuerdo con la invariancia (i) de la \cref{prop:invariancias}.
\end{observacion}

\section{Post-norm y pre-norm}

\begin{definicion}[Posición de la normalización]\label{def:pre-post}
Sea $\Norm$ una normalización (LayerNorm o RMSNorm) y $F$ una rama residual. Un bloque es
\[
\text{\emph{post-norm} si } h_\ell=\Norm\bigl(h_{\ell-1}+F(h_{\ell-1})\bigr),\qquad
\text{\emph{pre-norm} si } h_\ell=h_{\ell-1}+F\bigl(\Norm(h_{\ell-1})\bigr).
\]
\end{definicion}

El Transformer original~\cite{vaswani} es post-norm. GPT-2~\cite[§2.3]{gpt2} trasladó la normalización a la entrada de cada subbloque y añadió una normalización final tras el último bloque; LLaMA~\cite{llama} usa pre-norm con RMSNorm, y la mayoría de los modelos abiertos posteriores la han adoptado.

\begin{proposicion}[Jacobianos de ambas posiciones]\label{prop:jac-pre-post}
Sea $N_\ell$ el jacobiano de la normalización del bloque $\ell$ y $J_\ell$ el de la rama $F_\ell$, cada uno evaluado en su argumento. Entonces
\[
\text{post-norm: }\frac{\partial h_L}{\partial h_0}=\prod_{\ell=L}^{1}N_\ell\,(I+J_\ell),\qquad
\text{pre-norm: }\frac{\partial h_L}{\partial h_0}=\prod_{\ell=L}^{1}\bigl(I+J_\ell N_\ell\bigr),
\]
donde los productos se escriben con los índices en orden decreciente. Solo el desarrollo de pre-norm contiene el término $I$.
\end{proposicion}

En pre-norm, el flujo residual no pasa por ninguna normalización y conserva el camino identidad. A cambio, su norma crece con la profundidad (\cref{prop:crecimiento}), y hace falta una normalización final antes de la proyección a logits. Xiong \emph{et al.}~\cite{xiong} demuestran, con argumentos de campo medio, que en post-norm el gradiente esperado de los parámetros próximos a la salida es grande en la inicialización, lo que obliga a un calentamiento largo de la tasa de aprendizaje (módulo~7); en pre-norm los gradientes están acotados y el calentamiento puede acortarse o suprimirse.

\addtocounter{figure}{1}%

\begin{observacion}[Sobre los modelos comerciales]
Anthropic y Google no han hecho pública la normalización que usan Claude y Gemini. El código publicado de Grok-1~\cite{grok1} incluye capas RMSNorm.
\end{observacion}


% =====================================================================
\chapter{Dropout y recorte del gradiente}
\label{tema:4}
% =====================================================================

Este tema reúne dos técnicas que actúan sobre el entrenamiento sin cambiar la arquitectura. El dropout introduce ruido multiplicativo en las activaciones para regularizar; el recorte del gradiente limita la longitud de cada paso para que un gradiente anómalo no destruya lo aprendido. Se demuestran las propiedades que justifican ambas: el dropout invertido es insesgado y su varianza se calcula en forma cerrada; el recorte por norma global conserva la dirección del gradiente y acota el paso. Se sigue~\cite[§7.12]{gbc}, \cite[§5.6]{d2l}, \cite[§9.6.1]{bishopdl} y las notas de CS231n~\cite{cs231n} para el dropout, con el artículo de Srivastava \emph{et al.}~\cite{srivastava}; y \cite[§8.2.4 y §10.11]{gbc} y el artículo de Pascanu, Mikolov y Bengio~\cite{pascanu} para el recorte.

\section{Dropout}

\begin{definicion}[Dropout invertido]\label{def:dropout}
Sea $h\in\R^{d}$ un vector de activaciones y $p\in[0,1)$. Durante el entrenamiento, el \emph{dropout invertido} con probabilidad $p$ sustituye $h$ por
\[
\tilde h_i=\frac{m_i\,h_i}{1-p},\qquad m_i\sim\Bern(1-p)\ \text{independientes},\qquad i=1,\dots,d.
\]
En evaluación se usa $\tilde h=h$.
\end{definicion}

Cada unidad se apaga con probabilidad $p$ y las que sobreviven se amplifican por $1/(1-p)$. La variante original~\cite{srivastava} no amplificaba durante el entrenamiento y, en cambio, multiplicaba los pesos por $1-p$ en la evaluación; la variante invertida traslada esa corrección al entrenamiento y deja intacto el código de inferencia~\cite{cs231n}.

\begin{proposicion}[Esperanza y varianza]\label{prop:dropout}
Con la notación de la \cref{def:dropout}, para $h$ fijo,
\[
\E[\tilde h_i]=h_i,\qquad \Var(\tilde h_i)=\frac{p}{1-p}\,h_i^{2}.
\]
\end{proposicion}

\begin{corolario}[Dropout antes de una capa lineal]\label{cor:dropout-lineal}
Sea $y=w\T\tilde h$ con $w\in\R^{d}$ fijo. Entonces $\E[y]=w\T h$ y $\Var(y)=\frac{p}{1-p}\sum_{i=1}^{d}w_i^{2}h_i^{2}$.
\end{corolario}

El corolario muestra que, para una capa lineal, evaluar sin máscara da exactamente la salida media sobre todas las máscaras. Goodfellow \emph{et al.}~\cite[§7.12]{gbc} interpretan el dropout como el entrenamiento simultáneo de los $2^{d}$ subredes que se obtienen al eliminar unidades, con parámetros compartidos; la regla de evaluación sin máscara, llamada \emph{regla de escalado de pesos}, es exacta para modelos sin no linealidades ocultas, como el anterior, y solo aproximada en redes profundas no lineales.

\addtocounter{definicion}{1}%

\addtocounter{definicion}{1}%

\section{Recorte del gradiente por norma global}

\begin{definicion}[Recorte por norma global]\label{def:recorte}
Sean $g^{(1)},\dots,g^{(K)}$ los gradientes de los $K$ tensores de parámetros de un modelo y sea $g\in\R^{n}$ su concatenación. Dado un umbral $c>0$, el \emph{recorte por norma global} sustituye cada $g^{(k)}$ por
\[
\hat g^{(k)}=\alpha\,g^{(k)},\qquad \alpha=\min\Bigl(1,\ \frac{c}{\norm{g}_2}\Bigr),
\]
con el convenio $\alpha=1$ si $g=0$.
\end{definicion}

\begin{proposicion}[Propiedades del recorte]\label{prop:recorte}
Sea $\hat g=\alpha g$ el gradiente recortado.
\begin{enumerate}
  \item $\hat g$ tiene la dirección y el sentido de $g$, y $\norm{\hat g}_2=\min(\norm{g}_2,\,c)$.
  \item Un paso de descenso $\theta\leftarrow\theta-\eta\hat g$ cumple $\norm{\Delta\theta}_2\le\eta c$.
  \item Si $\hat g\ne 0$, $\hat g$ es una dirección de descenso: $\langle\nabla\Loss,-\hat g\rangle<0$ cuando $\hat g$ se calcula con el gradiente exacto.
  \item Recortar cada tensor por separado con el mismo umbral conserva, en general, la dirección de cada $g^{(k)}$ pero no la de $g$.
\end{enumerate}
\end{proposicion}

\addtocounter{definicion}{1}%

El recorte lo propusieron Pascanu, Mikolov y Bengio~\cite[§3.2]{pascanu} para redes recurrentes, donde el \cref{cor:pascanu} permite gradientes que explotan; en un acantilado de la superficie de pérdida~\cite[§8.2.4]{gbc}, el gradiente apunta en una dirección razonable pero tiene una longitud desproporcionada, y el recorte conserva lo primero y corrige lo segundo. En PyTorch se aplica con \texttt{torch.nn.utils.clip\_grad\_norm\_} entre \texttt{backward()} y el paso del optimizador.

\addtocounter{definicion}{1}%


% =====================================================================
\chapter{Aritmética de coma flotante y precisión mixta}
\label{tema:5}
% =====================================================================

Los modelos de lenguaje se entrenan con números de 16 e incluso de 8 bits. Un formato de coma flotante reparte sus bits entre el exponente, que fija el rango, y la mantisa, que fija la precisión relativa; esa elección determina qué valores desbordan, cuáles se anulan y qué sumandos se pierden al acumular. Este tema define los formatos con la norma IEEE~754~\cite{ieee754}, demuestra las cotas de rango y de error de redondeo, estudia la absorción de sumandos pequeños y justifica con ella los tres ingredientes de la precisión mixta de Micikevicius \emph{et al.}~\cite{micikevicius}. Se sigue además el tutorial de Goldberg~\cite{goldberg} y el libro de Higham~\cite{higham}.

\section{Formatos de coma flotante}

\begin{definicion}[Formato binario de coma flotante]\label{def:formato}
Sean $e\ge 2$ y $m\ge 1$ enteros. Un \emph{formato binario} con $e$ bits de exponente y $m$ bits de mantisa representa cada número con $1+e+m$ bits: un signo $s\in\{0,1\}$, un exponente almacenado $E\in\{0,\dots,2^{e}-1\}$ y una fracción $f=(f_1\dots f_m)_2$. Con el \emph{sesgo} $b=2^{e-1}-1$:
\begin{enumerate}
  \item si $1\le E\le 2^{e}-2$, el número es \emph{normal} y vale $x=(-1)^{s}\,(1.f_1\dots f_m)_2\cdot 2^{E-b}$;
  \item si $E=0$, vale $(-1)^{s}\,(0.f_1\dots f_m)_2\cdot 2^{1-b}$: el cero si $f=0$ y un número \emph{subnormal} si $f\ne 0$;
  \item si $E=2^{e}-1$, representa $\pm\infty$ si $f=0$ y \NaN\ (\emph{not a number}) si $f\ne 0$.
\end{enumerate}
\end{definicion}

Los formatos binary32 (fp32, $e=8$, $m=23$) y binary16 (fp16, $e=5$, $m=10$) son los de la norma~\cite{ieee754}. El formato bfloat16 (bf16, $e=8$, $m=7$) no forma parte de ella; se obtiene conservando los 16 bits más significativos de un fp32.

\begin{proposicion}[Rango y precisión]\label{prop:rango}
En el formato de la \cref{def:formato}:
\begin{enumerate}
  \item el mayor número finito es $x_{\max}=(2-2^{-m})\,2^{b}$;
  \item el menor normal positivo es $x_{\min}=2^{1-b}$, y el menor subnormal positivo, $2^{1-b-m}$;
  \item en cada intervalo $[2^{k},2^{k+1})$ con $1-b\le k\le b$, los números representables están equiespaciados con separación $2^{k-m}$;
  \item el \emph{épsilon de máquina}, distancia de~1 al siguiente número representable, es $\varepsilon=2^{-m}$.
\end{enumerate}
\end{proposicion}

\begin{table}[h]
\centering\small
\begin{tabular}{@{}lcccccc@{}}
\toprule
Formato & $e$ & $m$ & $x_{\max}$ & $x_{\min}$ (normal) & $\varepsilon=2^{-m}$ & Cifras decimales \\
\midrule
fp32 & 8 & 23 & $3{,}40\cdot 10^{38}$ & $1{,}18\cdot 10^{-38}$ & $1{,}19\cdot 10^{-7}$ & $\approx 7$ \\
fp16 & 5 & 10 & $65\,504$ & $6{,}10\cdot 10^{-5}$ & $9{,}77\cdot 10^{-4}$ & $\approx 3$ \\
bf16 & 8 & 7 & $3{,}39\cdot 10^{38}$ & $1{,}18\cdot 10^{-38}$ & $7{,}81\cdot 10^{-3}$ & $\approx 2$ \\
fp8 E4M3 & 4 & 3 & $448$ & $1{,}56\cdot 10^{-2}$ & $0{,}125$ & $\approx 1$ \\
fp8 E5M2 & 5 & 2 & $57\,344$ & $6{,}10\cdot 10^{-5}$ & $0{,}25$ & $<1$ \\
\bottomrule
\end{tabular}
\caption{Formatos usados en el entrenamiento de redes. El formato E4M3 no sigue la \cref{def:formato}(iii); véase la \cref{obs:e4m3}.}
\label{tab:formatos}
\end{table}

\addtocounter{definicion}{1}%

\begin{observacion}[El formato E4M3]\label{obs:e4m3}
Los dos formatos de 8 bits proceden de Micikevicius \emph{et al.}~\cite{fp8}. E5M2 sigue las reglas IEEE. E4M3 renuncia a los infinitos y reserva para \NaN\ un único patrón de bits (exponente y fracción con todos los bits a~1). Así, el exponente $E=15$ queda disponible para números normales y el máximo es $(1{,}110)_2\cdot 2^{15-7}=1{,}75\cdot 2^{8}=448$, en lugar del valor $1{,}875\cdot 2^{7}=240$ que darían las reglas IEEE.
\end{observacion}

El formato bf16 tiene el exponente de fp32 y, por tanto, prácticamente su mismo rango, con una precisión mucho menor. Pasar de fp32 a bf16 equivale a redondear la mantisa a 7 bits. Por su rango, bf16 es el formato habitual para entrenar modelos de lenguaje: casi nunca desborda.

\section{Redondeo y unidad de redondeo}

\begin{definicion}[Redondeo al más próximo]\label{def:fl}
Para $x\in\R$ con $\abs{x}\le x_{\max}$, $\fl(x)$ es el número representable más cercano a $x$; en caso de empate, el de último bit de mantisa par. La \emph{unidad de redondeo} es $u=\varepsilon/2=2^{-m-1}$.
\end{definicion}

\begin{teorema}[Error relativo de redondeo]\label{teo:redondeo}
Si $x_{\min}\le\abs{x}\le x_{\max}$, entonces $\fl(x)=x(1+\delta)$ con $\abs{\delta}\le u$.
\end{teorema}

La norma IEEE~754 exige que la suma, la resta, el producto, el cociente y la raíz cuadrada se redondeen correctamente: el resultado calculado es $\fl(x\circ y)$, el redondeo del resultado exacto. Junto con el \cref{teo:redondeo}, esto da el \emph{modelo estándar} del análisis de errores~\cite{higham,goldberg}: $\fl(x\circ y)=(x\circ y)(1+\delta)$ con $\abs\delta\le u$, mientras no haya desbordamiento ni subdesbordamiento.

\begin{lema}[Absorción]\label{lem:absorcion}
Sea $x>0$ representable con $2^{k}\le x<2^{k+1}$ y sea $0\le y<2^{k-m-1}$, la mitad de la separación del intervalo. Entonces $\fl(x+y)=x$.
\end{lema}

\begin{corolario}[Estancamiento de una suma]\label{cor:estancamiento}
Sea $a>0$ representable y sea la suma recursiva $s_0=0$, $s_{j+1}=\fl(s_j+a)$. Si para algún $j$ se cumple $2^{k}\le s_j<2^{k+1}$ con $a<2^{k-m-1}$, entonces $s_i=s_j$ para todo $i\ge j$. En particular, la suma no supera la menor potencia $2^{k}$ que cumple $2^{k}>a\,2^{m+1}$, más una separación.
\end{corolario}

\begin{table}[h]
\centering\small
\begin{tabular}{@{}lccc@{}}
\toprule
Operación & fp32 & fp16 & bf16 \\
\midrule
$300\cdot 300$ & $90\,000$ & \Inf & $90\,112$ \\
$e^{12}$ & $162\,754{,}8$ & \Inf & $162\,816$ \\
$1+2^{-8}$ & $1{,}0039$ & $1{,}0039$ & $1$ \\
$1000+1$ & $1001$ & $1001$ & $1000$ \\
$0{,}001$ sumado $10^{4}$ veces & $10{,}0004$ & $4$ & $0{,}5$ \\
\bottomrule
\end{tabular}
\caption{Resultados de operaciones sencillas en tres formatos, comprobados con PyTorch.}
\label{tab:errores}
\end{table}

\addtocounter{definicion}{1}%

\begin{teorema}[Error de la suma recursiva]\label{teo:suma}
Sea $\hat s_n$ el resultado de sumar $x_1,\dots,x_n$ en orden, $\hat s_1=x_1$, $\hat s_k=\fl(\hat s_{k-1}+x_k)$, sin desbordamiento ni subdesbordamiento, y sea $(n-1)u<1$. Entonces
\[
\abs{\hat s_n-\sum_{i=1}^{n}x_i}\le\gamma_{n-1}\sum_{i=1}^{n}\abs{x_i},\qquad\gamma_k=\frac{ku}{1-ku}.
\]
\end{teorema}

El teorema explica la diferencia entre formatos. En fp32, $u=2^{-24}=5{,}96\cdot 10^{-8}$ y, para $n=10^{4}$ sumandos de $0{,}001$, la cota es $\gamma_{9999}\cdot 10=0{,}006$; el error observado en la \cref{tab:errores} es $4\cdot 10^{-4}$. En fp16, $u=2^{-11}$ y $(n-1)u\approx 4{,}9>1$: la cota no garantiza nada, y la suma se estanca.

\section{Precisión mixta}

Micikevicius \emph{et al.}~\cite[§3]{micikevicius} entrenan en fp16 sin perder precisión con tres técnicas.

\paragraph{Copia maestra de los pesos en fp32 (§3.1).} El optimizador actualiza una copia fp32 de los pesos; la copia fp16 solo se usa en las pasadas hacia delante y hacia atrás. Por el \cref{lem:absorcion}, si un peso $w\in[2^{k},2^{k+1})$ recibe una actualización $\abs{\Delta w}<2^{k-m-1}$, la suma $w+\Delta w$ redondeada devuelve $w$. Como $w<2^{k+1}$, esto ocurre seguro si $\abs{w}/\abs{\Delta w}>2^{m+2}$ y puede ocurrir desde $2^{m+1}$; en fp16 este último cociente es $2^{11}=2048$, el que señala el artículo.

\paragraph{Escalado de la pérdida (§3.2).} Muchos gradientes de activaciones son menores que el menor subnormal de fp16, $2^{-24}$, y se anulan. Se multiplica la pérdida por un factor $S$, se retropropaga y se dividen los gradientes de los pesos por $S$ antes de la actualización, que se omite si aparece un \Inf\ o un \NaN. La elección $S=2^{j}$ es conveniente por la proposición siguiente.

\begin{proposicion}\label{prop:escala}
Si $x$ y $2^{j}x$ están en el rango normal, $\fl(2^{j}x)=2^{j}\fl(x)$. En particular, multiplicar o dividir un número representable por $2^{j}$ es exacto en ausencia de desbordamiento.
\end{proposicion}

\paragraph{Acumulación en fp32 (§3.3).} Los productos escalares de las multiplicaciones de matrices acumulan las sumas parciales en fp32 y solo convierten el resultado a fp16; las reducciones grandes (sumas de la normalización y de la softmax) también se hacen en fp32. El \cref{teo:suma} justifica esta práctica.

\begin{observacion}[bf16 y fp8]
Con bf16 no hace falta escalar la pérdida, porque su rango es el de fp32, pero la copia maestra y la acumulación en fp32 siguen siendo necesarias, porque su precisión es aún menor que la de fp16. Los formatos fp8 se usan en las GPU recientes para las multiplicaciones de matrices, E4M3 en la pasada hacia delante y E5M2 para los gradientes, siempre con factores de escala por tensor~\cite{fp8}; el módulo~17 los estudia con el hardware y el módulo~19 trata la cuantización para inferencia. Anthropic, Google y xAI no han hecho pública la precisión con que entrenan Claude, Gemini y Grok.
\end{observacion}


% =====================================================================
\chapter{Estabilidad numérica de softmax y entropía cruzada}
\label{tema:6}
% =====================================================================

La capa de salida de un modelo de lenguaje calcula una softmax sobre decenas de miles de logits, y la pérdida de entrenamiento es su entropía cruzada (módulo~1). Ambas contienen exponenciales, que en coma flotante desbordan o se anulan con argumentos moderados. Este tema determina esos límites, introduce la función log-sum-exp y su evaluación estable, deduce de ella la forma correcta de calcular la entropía cruzada a partir de los logits y termina con la propagación de \NaN\ y un procedimiento para diagnosticar una pérdida que deja de ser finita. Se sigue~\cite[§4.1]{gbc}, \cite[§4.5.2]{d2l} y las prácticas de CS336~\cite{cs336}.

\section{Límites de la exponencial}

\begin{proposicion}\label{prop:exp-rango}
En un formato con mayor número finito $x_{\max}$ y menor normal $x_{\min}$, el valor exacto $e^{t}$ es finito y normal si y solo si $\ln x_{\min}\le t\le\ln x_{\max}$. Los intervalos son, aproximadamente,
\[
\text{fp32 y bf16: }[-87{,}3;\ 88{,}7],\qquad \text{fp16: }[-9{,}70;\ 11{,}09].
\]
\end{proposicion}

Por debajo del límite inferior, el resultado es subnormal, con menos bits de precisión, y algo más abajo se anula. Goodfellow \emph{et al.}~\cite[§4.1]{gbc} describen las dos formas en que falla la softmax ingenua $e^{z_i}/\sum_je^{z_j}$: si algún logit es grande, el numerador y el denominador desbordan y el cociente $\infty/\infty$ es \NaN; si todos son muy negativos, ambos se anulan y $0/0$ es \NaN.

\section{La función log-sum-exp}

\begin{definicion}[Log-sum-exp]\label{def:lse}
Para $z\in\R^{n}$, $\LSE(z)=\ln\sum_{i=1}^{n}e^{z_i}$. La softmax del módulo~1 se escribe $\softmax(z)_i=e^{z_i-\LSE(z)}$.
\end{definicion}

\begin{proposicion}[Propiedades de log-sum-exp]\label{prop:lse}
Sea $z\in\R^{n}$ y $c\in\R$.
\begin{enumerate}
  \item $\LSE(z+c\uno)=\LSE(z)+c$ y $\softmax(z+c\uno)=\softmax(z)$.
  \item $\max_iz_i\le\LSE(z)\le\max_iz_i+\ln n$; la cota superior se alcanza si y solo si todas las componentes son iguales.
  \item $\nabla\LSE(z)=\softmax(z)$.
  \item $\nabla^{2}\LSE(z)=\diag(p)-pp\T$ con $p=\softmax(z)$, que es semidefinida positiva; en consecuencia, $\LSE$ es convexa.
\end{enumerate}
\end{proposicion}

\begin{proposicion}[Evaluación estable]\label{prop:lse-estable}
Sea $c=\max_iz_i$. Entonces
\[
\LSE(z)=c+\ln S,\qquad \softmax(z)_i=\frac{e^{z_i-c}}{S},\qquad S=\sum_{i=1}^{n}e^{z_i-c}\in[1,n].
\]
En esta forma ninguna exponencial desborda, porque todos sus argumentos son $\le 0$, y el denominador nunca se anula, porque contiene el término $e^{0}=1$.
\end{proposicion}

\addtocounter{definicion}{1}%

\section{Entropía cruzada a partir de los logits}

\begin{proposicion}[Log-softmax y entropía cruzada]\label{prop:ce}
Sea $z\in\R^{n}$ el vector de logits y $y\in\{1,\dots,n\}$ la clase correcta. Entonces
\[
\ln\softmax(z)_i=z_i-\LSE(z),\qquad \ell(z,y)=-\ln\softmax(z)_y=\LSE(z)-z_y\ge 0,
\]
\[
\nabla_z\ell=\softmax(z)-e_y,
\]
donde $e_y$ es el vector canónico.
\end{proposicion}

La fórmula $\ell=\LSE(z)-z_y$, calculada con la \cref{prop:lse-estable}, no evalúa nunca el logaritmo de un número que pueda haberse anulado. Por eso las funciones de pérdida de las bibliotecas, como \texttt{F.cross\_entropy} de PyTorch, reciben logits y no probabilidades~\cite[§4.5.2]{d2l}, y la composición \texttt{torch.log(torch.softmax(z))} debe evitarse.

\addtocounter{definicion}{1}%

\section{Propagación de \texorpdfstring{\NaN}{NaN} y diagnóstico}

La norma IEEE~754~\cite{ieee754} define como \NaN\ los resultados de las operaciones inválidas ${\infty-\infty}$, ${0\cdot\infty}$, ${0/0}$, ${\infty/\infty}$ y $\sqrt{x}$ con $x<0$, y establece que toda operación aritmética con un operando \NaN\ da \NaN; además, toda comparación con \NaN\ es falsa. En particular, $\NaN\cdot 0=\NaN$: multiplicar por cero no limpia el valor.

\begin{observacion}[Contagio]\label{obs:contagio}
Si la pérdida de un paso es \NaN, la retropropagación multiplica ese valor por todos los jacobianos locales, de modo que el gradiente de todo parámetro del que depende la pérdida es \NaN, y la actualización convierte esos parámetros en \NaN. A partir de ese paso la red ya no produce números. Por eso el diagnóstico debe hacerse en el primer paso en que aparece un valor no finito.
\end{observacion}

\begin{table}[h]
\centering\small
\begin{tabularx}{\textwidth}{@{}>{\raggedright\arraybackslash}X>{\raggedright\arraybackslash}X>{\raggedright\arraybackslash}X@{}}
\toprule
Síntoma & Causa probable & Remedio \\
\midrule
La pérdida crece varios órdenes de magnitud y luego da \NaN & tasa de aprendizaje alta; gradientes que explotan (\cref{tema:1}) & bajar la tasa, calentamiento, recorte del gradiente (\cref{tema:4}) \\
\Inf\ repentino en fp16 & desbordamiento por encima de $65\,504$ & bf16, o escalado dinámico de la pérdida (\cref{tema:5}) \\
\NaN\ en una softmax o un logaritmo & exponencial desbordada, $\ln 0$ & \cref{prop:lse-estable}; entropía cruzada con logits \\
\NaN\ en una normalización & división por $\RMS(x)=0$ & $\epsilon>0$ dentro de la raíz (\cref{tema:3}) \\
Picos aislados de la pérdida & lote anómalo & recorte, descartar el lote, volver a un punto de control \\
\bottomrule
\end{tabularx}
\caption{Diagnóstico de una pérdida no finita.}
\label{tab:diagnostico}
\end{table}

\addtocounter{figure}{1}%

\begin{observacion}[Lo que se observa en la práctica]
En la práctica del módulo se provoca el fallo a propósito: una red residual sin normalizar entrenada con tasa de aprendizaje~10 pasa de una pérdida de $2{,}6$ a $10^{30}$ y llega a \NaN\ en el paso~24. Con recorte del gradiente a norma~1, la pérdida se mantiene finita durante los 60 pasos.
\end{observacion}



\backmatter
% =====================================================================
\begin{thebibliography}{99}
\addcontentsline{toc}{chapter}{Bibliografía}
\markboth{Bibliografía}{}
\small
\def\UrlBreaks{\do\/\do-}

\bibitem{layernorm}
J.~L. Ba, J.~R. Kiros y G.~E. Hinton.
\newblock Layer normalization.
\newblock Prepublicación, 2016.
\newblock \url{https://arxiv.org/abs/1607.06450}

\bibitem{bishopdl}
C.~M. Bishop y H.~Bishop.
\newblock \emph{Deep Learning: Foundations and Concepts}.
\newblock Springer, 2024.
\newblock Cap.~7, «Gradient Descent» (§7.4 «Normalization»); cap.~9, «Regularization» (§9.5 «Residual Connections» y §9.6.1 «Dropout»).
\newblock \url{https://www.bishopbook.com/}

\bibitem{cs231n}
F.-F. Li \emph{et al.}
\newblock CS231n: Deep Learning for Computer Vision. Notas del curso.
\newblock Stanford University, 2026. «Neural Networks Part 2: Setting up the Data and the Loss», apartados «Weight Initialization» y «Regularization».
\newblock \url{https://cs231n.github.io/neural-networks-2/}

\bibitem{cs224n}
C.~Manning \emph{et al.}
\newblock CS224N: Natural Language Processing with Deep Learning.
\newblock Stanford University, invierno de 2021. Práctica~3, «Dependency Parsing».
\newblock \url{https://web.stanford.edu/class/archive/cs/cs224n/cs224n.1214/assignments/a3.pdf}

\bibitem{cs336}
T.~Hashimoto y P.~Liang.
\newblock CS336: Language Modeling from Scratch.
\newblock Stanford University, primavera de 2026. Práctica~1, «Basics», versión~26.0.3.
\newblock \url{https://github.com/stanford-cs336/assignment1-basics}

\bibitem{cs336b}
T.~Hashimoto y P.~Liang.
\newblock CS336: Language Modeling from Scratch.
\newblock Stanford University, primavera de 2026. Práctica~2, «Systems», versión~26.1.3.
\newblock \url{https://github.com/stanford-cs336/assignment2-systems}

\bibitem{glorot}
X.~Glorot e Y.~Bengio.
\newblock Understanding the difficulty of training deep feedforward neural networks.
\newblock En \emph{Proceedings of the 13th International Conference on Artificial Intelligence and Statistics (AISTATS)}, PMLR 9:249--256, 2010.

\bibitem{goldberg}
D.~Goldberg.
\newblock What every computer scientist should know about floating-point arithmetic.
\newblock \emph{ACM Computing Surveys}, 23(1):5--48, 1991.

\bibitem{gbc}
I.~Goodfellow, Y.~Bengio y A.~Courville.
\newblock \emph{Deep Learning}.
\newblock MIT Press, 2016.
\newblock Cap.~4, «Numerical Computation»; cap.~7, «Regularization for Deep Learning»; cap.~8, «Optimization for Training Deep Models»; cap.~10, «Sequence Modeling: Recurrent and Recursive Nets».
\newblock \url{https://www.deeplearningbook.org/}

\bibitem{resnet}
K.~He, X.~Zhang, S.~Ren y J.~Sun.
\newblock Deep residual learning for image recognition.
\newblock En \emph{IEEE Conference on Computer Vision and Pattern Recognition (CVPR)}, págs.~770--778, 2016.
\newblock \url{https://arxiv.org/abs/1512.03385}

\bibitem{he-init}
K.~He, X.~Zhang, S.~Ren y J.~Sun.
\newblock Delving deep into rectifiers: Surpassing human-level performance on ImageNet classification.
\newblock En \emph{IEEE International Conference on Computer Vision (ICCV)}, 2015.
\newblock \url{https://arxiv.org/abs/1502.01852}

\bibitem{resnet-id}
K.~He, X.~Zhang, S.~Ren y J.~Sun.
\newblock Identity mappings in deep residual networks.
\newblock En \emph{European Conference on Computer Vision (ECCV)}, 2016.
\newblock \url{https://arxiv.org/abs/1603.05027}

\bibitem{higham}
N.~J. Higham.
\newblock \emph{Accuracy and Stability of Numerical Algorithms}, 2.ª ed.
\newblock SIAM, 2002.

\bibitem{ieee754}
IEEE Computer Society.
\newblock \emph{IEEE Standard for Floating-Point Arithmetic}, IEEE Std 754-2019.
\newblock IEEE, 2019.

\bibitem{ioffe}
S.~Ioffe y C.~Szegedy.
\newblock Batch normalization: Accelerating deep network training by reducing internal covariate shift.
\newblock En \emph{International Conference on Machine Learning (ICML)}, 2015.
\newblock \url{https://arxiv.org/abs/1502.03167}

\bibitem{micikevicius}
P.~Micikevicius, S.~Narang, J.~Alben, G.~Diamos, E.~Elsen, D.~Garcia, B.~Ginsburg, M.~Houston, O.~Kuchaiev, G.~Venkatesh y H.~Wu.
\newblock Mixed precision training.
\newblock En \emph{International Conference on Learning Representations (ICLR)}, 2018.
\newblock \url{https://arxiv.org/abs/1710.03740}

\bibitem{fp8}
P.~Micikevicius \emph{et al.}
\newblock FP8 formats for deep learning.
\newblock Prepublicación, 2022.
\newblock \url{https://arxiv.org/abs/2209.05433}

\bibitem{pascanu}
R.~Pascanu, T.~Mikolov e Y.~Bengio.
\newblock On the difficulty of training recurrent neural networks.
\newblock En \emph{International Conference on Machine Learning (ICML)}, PMLR 28(3):1310--1318, 2013.
\newblock \url{https://arxiv.org/abs/1211.5063}

\bibitem{gpt2}
A.~Radford, J.~Wu, R.~Child, D.~Luan, D.~Amodei e I.~Sutskever.
\newblock Language models are unsupervised multitask learners.
\newblock Informe técnico, OpenAI, 2019.

\bibitem{cmu11785}
B.~Raj \emph{et al.}
\newblock 11-785: Introduction to Deep Learning.
\newblock Carnegie Mellon University, otoño de 2025.
\newblock \url{https://deeplearning.cs.cmu.edu/}

\bibitem{srivastava}
N.~Srivastava, G.~Hinton, A.~Krizhevsky, I.~Sutskever y R.~Salakhutdinov.
\newblock Dropout: A simple way to prevent neural networks from overfitting.
\newblock \emph{Journal of Machine Learning Research}, 15(56):1929--1958, 2014.
\newblock \url{https://jmlr.org/papers/v15/srivastava14a.html}

\bibitem{llama}
H.~Touvron \emph{et al.}
\newblock LLaMA: Open and efficient foundation language models.
\newblock Prepublicación, 2023.
\newblock \url{https://arxiv.org/abs/2302.13971}

\bibitem{vaswani}
A.~Vaswani, N.~Shazeer, N.~Parmar, J.~Uszkoreit, L.~Jones, A.~N. Gomez, Ł.~Kaiser e I.~Polosukhin.
\newblock Attention is all you need.
\newblock En \emph{Advances in Neural Information Processing Systems 30 (NeurIPS)}, 2017.
\newblock \url{https://arxiv.org/abs/1706.03762}

\bibitem{grok1}
xAI.
\newblock Grok-1: código y pesos publicados.
\newblock 2024.
\newblock \url{https://github.com/xai-org/grok-1}

\bibitem{xiong}
R.~Xiong, Y.~Yang, D.~He, K.~Zheng, S.~Zheng, C.~Xing, H.~Zhang, Y.~Lan, L.~Wang y T.-Y. Liu.
\newblock On layer normalization in the Transformer architecture.
\newblock En \emph{International Conference on Machine Learning (ICML)}, 2020.
\newblock \url{https://arxiv.org/abs/2002.04745}

\bibitem{d2l}
A.~Zhang, Z.~C. Lipton, M.~Li y A.~J. Smola.
\newblock \emph{Dive into Deep Learning}, versión 1.0.3.
\newblock Cambridge University Press, 2023.
\newblock Cap.~4, «Linear Neural Networks for Classification»; cap.~5, «Multilayer Perceptrons»; cap.~8, «Modern Convolutional Neural Networks»; cap.~9, «Recurrent Neural Networks».
\newblock \url{https://d2l.ai/}

\bibitem{rmsnorm}
B.~Zhang y R.~Sennrich.
\newblock Root mean square layer normalization.
\newblock En \emph{Advances in Neural Information Processing Systems 32 (NeurIPS)}, 2019.
\newblock \url{https://arxiv.org/abs/1910.07467}

\end{thebibliography}


\end{document}
